%----------------------------------------------------------------------------
\chapter{Code, tools and checks}\label{app:code}
%----------------------------------------------------------------------------
\section{GitHub}\label{app:code:github}
Keep your work in the three GitHub repositories described in \cref{app:start:access}: code, models, notes, drawings, PDFs, slides and the Overleaf-synced paper. Readers can access code anonymously for double-blind review through \href{https://anonymous.4open.science/}{Anonymous GitHub}, or as a permanent DOI archive through Zenodo (\cref{app:zenodo}).

\paragraph{Repository layout.} The notes repository's submodules (\cref{app:start:access}) keep the histories separate within one clone:
\begin{itemize}[noitemsep]
  \item notes repository: \texttt{README.md} linking the pieces; \texttt{notes/} for notes and drawings, with consultation summaries in \texttt{notes/consultations/}; \texttt{slides/} for consultation slides and the presentation;
  \begin{itemize}[noitemsep]
    \item \texttt{paper/}: submodule, the Overleaf-synced paper repository, with the TikZ sources of every figure in \texttt{figures/};
    \item \texttt{code/}: submodule, the code repository: the code in \texttt{src/} and \texttt{tests/}, the measurement scripts and their raw results in \texttt{experiments/}, and at the root \texttt{README.md}, \texttt{LICENSE}, \texttt{.zenodo.json}, \texttt{CITATION.cff}, the \texttt{justfile}, the \texttt{Dockerfile} and \texttt{.devcontainer/}.
  \end{itemize}
\end{itemize}
Include the code and data behind the paper's results in each release (\cref{app:zenodo}).

\paragraph{One command.} The \href{https://github.com/BlackLight54/tdk-code-template}{tdk-code-template} provides the layout above, Zenodo files, Docker images and one GitHub Actions workflow per recipe: \texttt{ci} (build, test and smoke reproduction), \texttt{evaluate}, \texttt{benchmark}, \texttt{robustness}, \texttt{demo} and \texttt{release}.

Put recurring commands in a \texttt{justfile} for the \href{https://github.com/casey/just}{just} command runner. We require that \texttt{just reproduce} regenerates every paper figure and table. Name individual recipes after their outputs: \texttt{just fig-scalability} for the scalability figure (\cref{fig:guide:scalability}) and \texttt{just tab-robustness} for the robustness table (\cref{tab:guide:robustness}). Provide \texttt{just reproduce --smoke} with a reduced workload that finishes in minutes. GitHub Actions runs it on every push to catch broken README commands.

\section{The checks and paper CI}\label{app:code:paperci}
The paper repository's \texttt{.github/workflows/paper.yml} builds both PDFs on every push to \texttt{main} and every pull request that is not a draft. It runs each rule in \texttt{scripts/lint.lua} as a separate step. The run page shows each result and annotates findings at the relevant file and line. The job summary starts with the essentials (\cref{app:writing:essentials}): a line of marks, then a row per essential. A second table groups the other rules as on the Checks page. There are three rule families; \texttt{texlua scripts/lint.lua list} describes each rule, with its severity and phase.
\begin{description}[noitemsep]
  \item[source] Checks the sources and build log. Gates cover phase-appropriate essentials (\texttt{essentials}), the template structure (\texttt{frame}) and a clean build (\texttt{build}, \texttt{overfull}); the chapter-order check (\texttt{spine}) only advises. They also check floats, captions, axes, tables, labels, references, acronyms, listings, the shape of the reference list, typography, consistent terms, that no TDK paper of the same year is cited or linked (\cref{app:tdk:companion}), headline numbers against the evaluation, and completion of the AI declaration and title page. The burn-down rules for placeholders and beats only report, and the \texttt{burndown} rule writes their counts as one line of JSON for CI to collect over time. A rule that cannot run, because a tool is missing or the profile excludes it, says \emph{skipped}, never \emph{pass}.
  \item[pdf] Checks the reviewer PDF against the form requirements (\cref{app:tdk:form}), even if you replaced the template. Checks cover the paper size, metadata, embedded fonts, a 15\,MB limit and the body font: \texttt{pdf-font} fails on a font that is not a serif or not 12\,pt, and only warns on a serif other than Times. The 15\,MB limit is our safeguard, taken from the faculty's thesis guideline (\texttt{include/guideline.tex}); it is not a TDK rule. They measure line spacing from line pitch and margins from the text on every page. They report text below 8\,pt, count main-text pages from the PDF outline, and check required parts, including both abstracts and the AI declaration. For a thesis, \texttt{LINT\_PROFILE=thesis} applies the faculty's thesis thresholds.
  \item[tools] Reports from TeX Live and container tools: chktex and lacheck for typography, texcount for chapter word counts and abstract length, checkcites for uncited reading-list entries, the BibTeX log, aspell on the detexed chapters for unknown words and mixed British and American spellings, and a link check. These report and never block the build, with one exception: \texttt{repo-hygiene} fails at every phase when the repository tracks build output such as a PDF, a log or an \texttt{.aux} file.
\end{description}
Each rule has a \emph{severity} and a \emph{phase}. A \emph{check} fails the run, \emph{advice} warns, and a \emph{count} reports numbers. A check whose phase is \texttt{always} fails at every phase: \texttt{essentials}, \texttt{frame}, \texttt{chapters}, \texttt{build}, \texttt{overfull}, \texttt{figures}, \texttt{labels}, \texttt{crefs}, \texttt{bib}, \texttt{companion}, \texttt{code}, \texttt{preamble} and \texttt{repo-hygiene}. The other checks become required at \texttt{draft}, \texttt{polish} or \texttt{submit}; before then, they only report. The phases are those of the lifecycle (\cref{app:lifecycle}): \texttt{topic}, \texttt{setup}, \texttt{research}, \texttt{plan}, \texttt{draft}, \texttt{polish} or \texttt{submit}. The first four share the checks' first level, and the lint rejects any other value. Page count blocks only on \texttt{main}, so a branch may be short while you draft a chapter. Locally, \texttt{make check} runs every rule against the last build, \texttt{make check-reviewer} against the reviewer PDF, and \texttt{make check-submit} against the reviewer PDF at the \texttt{submit} phase whatever \texttt{\textbackslash lintphase} says, the check before you upload; run one rule with, for example, \texttt{make lint-essentials}. The same lint generates the draft PDF's \emph{Checks} page on every compile (\cref{app:latex:todo}), keeping the PDF and CI consistent.

The workflow is part of the Overleaf template and is copied into each new student project. Overleaf's GitHub sync creates the student's repository; its first push starts CI. View runs under GitHub's \emph{Actions} tab or on a pull request.
\paragraph{Lifecycle CI.} The notes repository from \href{https://github.com/BlackLight54/tdk-notes-template}{tdk-notes-template} has a separate workflow. Every push and Monday morning, it regenerates \texttt{lifecycle.md} with file-based checks. It checks that the three anchors in \texttt{notes/anchors/} are longer than stubs, counts summaries in \texttt{notes/consultations/}, and checks that submodules have moved within two weeks. It also checks for the code README, license, justfile, tests, Dockerfile and Zenodo file, and requires slides before the conference. Due phases come from the paper's \texttt{\textbackslash lintphase}.

Paper CI uploads the paper's burn-down counts as \texttt{burndown.json} with every run; the notes repository's lifecycle workflow appends them to \texttt{burndown.csv}, plots \texttt{burndown.svg} and commits both for consultation use. These checks inventory the work; we still assess its quality.

\advisortodo{Move \texttt{tdk-notes-template} and \texttt{tdk-code-template} to an FTSRG organisation when there is one for the lab; \texttt{ftsrg-bta} is the course's.}

\section{Running the reviewer passes}\label{app:code:reviewers}
\paragraph{Self-score and reviewer passes.} The paper repository carries two agent skills in \path{.claude/skills/}, so each commit has the skills that review it. On every push, paper CI uploads the evidence of that commit as one artifact: both PDFs, the build logs, the lint output and the skim view, kept for 14 days. \path{.claude/skills/self-score/} scores the artifact's \texttt{reviewer.pdf} A--E using the review sheet's wording and page references, and identifies the weakest criterion and cheapest fix. \path{.claude/skills/reviewer-passes/} runs the skim reviewer and expert panel on the same artifact. \texttt{scripts/review/review.py} fetches the artifact of a named commit and checks its hashes. Self-score writes nothing; the reviewer passes write nothing unless you record a run. \Cref{app:lifecycle:reviewers} says what the passes are and when to run them; this section says how.

\paragraph{Running them.} Push your commit and wait for paper CI, which builds the evidence of every push (\cref{app:code:paperci}). Then, in your paper repository, ask your agent to run the reviewer passes, or name one persona: \enquote{review as the cryptographer}. The skill in \path{.claude/skills/reviewer-passes/} fetches the CI artifact of that commit with \texttt{scripts/review/review.py}, which refuses an artifact whose hashes do not match. It gives each persona a directory holding only that persona's inputs: the reviewer PDF, the skim view, or both, and your spine if you name the file. Without a spine, a persona judges overclaiming against the abstract and the introduction. Each persona runs apart from the others and never sees another persona's findings or an earlier run. The chat response gives each persona's worst finding, the criterion and page, and the cheapest fix. It then lists the three cheapest changes overall. Ask a persona for more to see its remaining findings.

\paragraph{Recording a run.} By default the skill writes no files and never edits the paper. You choose the fixes, and the consultation summary (\cref{app:consult}) records your decisions. We record a run at three milestones (the spine premortem, the full draft and before upload), and you may record any run you choose. A recorded run becomes one issue on your paper repository: the panel, the skim scores, one checkbox per finding, and the hashes of the artifact and of each persona's inputs. Propose a disposition for each finding in a comment on the issue: \texttt{F3 fixed in <sha>}, \texttt{F4 rejected: <reason>} or \texttt{F5 deferred}. We confirm the dispositions with the label \texttt{review-confirmed}. The issue is not an input to the next run, and the skill keeps no memory between runs. Rerun it on the revised paper to check whether the fixes answer the objections.

Ignore findings that do not identify a place in the paper. Verify every suggested source before citing it; a persona may invent one. Tell us if two personas raise the same concern or one argues outside its remit, so we can fix their boundaries in the persona files.

\section{The README}\label{app:code:readme}
Every TDK needs a \emph{results} repository that lets reviewers reproduce its numbers. Some also produce a reusable \emph{tool}. The code template's \texttt{README.md} starts with the results sections. Extra tool sections appear in a commented block below: delete the block for results only, or uncomment it for a tool.

\paragraph{The results README.} In this order:
\begin{enumerate}[noitemsep]
  \item What this is, in two sentences, with the paper's title and the concept DOI badge (\cref{app:zenodo}).
  \item The reproduction command: \texttt{docker run \dots{} reproduce} for every figure and table (\cref{app:code:docker}). Follow it with \texttt{just reproduce} for local clones and a list of pinned dependencies. Keep dependencies to a minimum.
  \item A results map with one row per paper figure and table: its number, reproduction command, output path and wall-clock time on specified hardware.
  \item Reproduction limits: secrets, special hardware, long runs or nondeterminism, including the seed policy and expected variance. Keep this section even if everything reproduces with fixed seeds.
  \item The results tag: which Git tag and Zenodo version the paper's numbers came from (\cref{app:zenodo:versions}).
  \item How to cite, and the license.
\end{enumerate}

\paragraph{The tool README} adds installation, use on the reader's own input with a worked example, contribution instructions and maintainers after the results sections. Follow the conventions of the open-source foundations we work with, such as \href{https://github.com/LF-Decentralized-Trust/lab-proposals/blob/main/README.md}{LF Decentralized Trust Labs}: an Apache 2.0 \texttt{LICENSE}, a \href{https://developercertificate.org/}{Developer Certificate of Origin} sign-off on every commit (\texttt{git commit -s}), \texttt{MAINTAINERS.md} and a \texttt{SECURITY.md} with vulnerability-reporting instructions. A compliant tool can be proposed as a lab after the TDK.

\section{Docker}\label{app:code:docker}
Make \texttt{docker run} the simplest way to reproduce results. We assume a computer science reviewer has Docker; see Docker's \href{https://docs.docker.com/get-started/}{Get started} if you are new to it. The code template provides two images:
\begin{itemize}[noitemsep]
  \item The \textbf{runtime image} is built by the root \texttt{Dockerfile}, with \texttt{just} as its entry point. Run \texttt{docker run --rm -v \$PWD/out:/out ghcr.io/<you>/<repo> reproduce} to regenerate all results into \texttt{out/}, or \texttt{\dots{} fig-scalability} for one figure. Pin every dependency to make the environment independent of the reviewer's machine.
  \item The \textbf{development image} includes the toolchain, language servers and test runner. Its \texttt{.devcontainer/devcontainer.json} follows the open \href{https://containers.dev/}{Development Containers} specification. VS Code, GitHub Codespaces and Docker's development environment feature can open the repository ready for work. We use the same environment when helping you.
\end{itemize}
A \href{https://docs.github.com/en/actions}{GitHub Actions} workflow builds both images on every push and publishes the runtime image to the \href{https://docs.github.com/en/packages/working-with-a-github-packages-registry/working-with-the-container-registry}{GitHub Container Registry} on each release. Use the same tag for the image and Zenodo version. We require the README's first command to select the results-tag image.

\section{Demos}\label{app:code:demo}
Record demonstrations with a script in the code repository, never by hand. A script lets you update the recording after code changes and rerun it before the presentation. Have CI publish a fresh recording on every push for the README, the presentation and reviewers.
\begin{itemize}[noitemsep]
  \item For a web UI, we recommend \href{https://playwright.dev/}{Playwright}: it drives the real workflow and saves a video and trace. Draw captions into the frames at each step.
  \item For a backend, CLI or prover such as Tamarin, Lean or a model checker, we recommend recording a scripted terminal session with \href{https://asciinema.org/}{asciinema}. Render it to SVG or GIF, or embed the cast beside the verdicts in HTML. We use this for our backends and proof assistants.
\end{itemize}
Use captions to explain the invariant, rejected attack or result in the running example's terms. Do more than narrate clicks. The code template includes a recorder in \texttt{demo/} and a \texttt{demo} workflow that uploads its recording as an artifact. Run it locally with \texttt{just demo}.

\section{Behaviour-driven tests}
For behaviour-driven development (BDD), you can write the expected behaviour as scenarios with \href{https://cucumber.io/}{Cucumber}; the scenarios then double as evaluation cases that a reviewer can read.

\section{AI tools}\label{app:code:ai}
\paragraph{Tools, as of October 2026.} These recommendations can become outdated within months.
\begin{itemize}
  \item For research, we recommend the strongest models. We prefer Claude Fable and ChatGPT Astra; ask whether we can help with a subscription.
  \item Prefer an agent working in your repository: Claude Code or Codex, ideally through T3 Code, GitKraken Kepler or herdr. This gives the agent access to your files and its repository tools. You \emph{need} version control when AI edits files.
  \item Plugins and skills: the Matt Pocock and poetato plugins for productivity; caveman, ponytail, squeez and TRON if you want to spare tokens. Three we use daily: \href{https://www.aihero.dev/skills-grill-me}{grill-me} (\href{https://claudemarketplaces.com/skills/mattpocock/skills/grill-me}{marketplace page}) interrogates a plan before you build it; \href{https://github.com/cursor/plugins/blob/main/pstack/skills/unslop/SKILL.md}{unslop} takes the machine out of generated prose; \href{https://github.com/DietrichGebert/ponytail/blob/main/skills/ponytail/SKILL.md}{ponytail} cuts the tokens an agent spends.
  \item We recommend concise instructions. Overloaded \texttt{AGENTS.md} files and memories can impair model performance. Prefer skills that clearly describe the range of tasks they apply to.
\end{itemize}
